% ============ 第 9 讲 移民与外包 ============
\chap{第 9 讲　移民与外包}{Migration \& Offshoring}

\textbf{本讲导读}：本讲对应 Lecture 9（Migration \& Offshoring, Oct.\ 23），共 66 页（其中 4 页整页板书截图）。内容：单部门移民模型（迁移均衡、移民与移民输出国的总福利与分配分析、摩擦与长期调整、技能结构）；政策讨论（边境墙）；外包（offshoring）的经济模型（价值链与``微笑曲线''、任务连续统、外包决策的阈值 $s^*$、贸易成本下降的比较静态、技能密集度与技能溢价、GVC 整合的含义）。\textbf{注}：文中页码为幻灯片自身编号；整页板书截图页无编号，标注为``板书''。

\secn{9.1 回顾与新的问题}

\textbf{行政通知（第 2 页）}：习题集 2 截止\textbf{11 月 1 日}（老师手写将打印版 Nov.\ 29 划掉改为 Nov.\ 1）；今天之前的内容全部期中（\textbf{Nov.\ 6}）可考；10 月 28 日为经验课（不考试）；\kw{11 月 4 日复习课}（前半复习理论、后半讨论 2025 春秋期中精选题，投票周二晚截止）。

\textbf{上一讲回顾（第 3--5 页）}：\term{长期（H-O 与 Rybczynski）}——两商品两要素的 H-O 中某要素增加：\kw{要素价格不敏感性（Factor price insensitivity，粉底高亮）}（$w/r$ 不变）；密集使用该要素的行业产出（与要素使用）增加、另一行业减少；关键假设：开放贸易、\kw{最终品价格不变}（红字）、\kw{多样化锥}（粉底高亮；禀赋比的范围）。\term{短期（特定要素）}——关键假设：最终品价格不变；\kw{移民}使两部门就业与生产扩张，\redkw{(国内)工人福利下降、特定要素所有者福利上升}；\kw{单个部门的 FDI} 使该部门扩张另一部门收缩，\redkw{工人福利上升、特定要素所有者福利下降}（板书手绘两张劳动市场图示）。

\textbf{今天的问题（第 6--7 页）}：\kw{单部门移民模型}（粉底高亮）——总福利与分配福利分析；\term{外包的经济模型}。移民决策与福利：上一讲：给定劳动流入，行业如何调整？今天：\kw{劳动为什么会移动、何时停止？}然后：\term{原有居民与移民的得与失}。单部门：撇开行业配置来研究移民——单一商品下固定价格不能同时钉住 \kw{w} 与 \kw{r}；一般得/失预测与短期 SFM 类似（手写 similar to SFM）。

\secn{9.2 单部门移民模型}

\begin{assumpbox}[模型设定（第 8 页）]
两国用\term{劳动 L 与固定资本 K} 生产一种商品（手写：$Y_i = F(K_i, L_i)$，$K_i$ \kw{specific to country}（国家特定）、$L_i$ \kw{may migrate}（可能移民））；技术相同、劳动同质、共同价格 $P = 1$。$L$ 增加时 $w/P$ 与 $r/P$ 如何变化？
\begin{itemize}
  \item \kw{$w/P = MPL$}（粉底高亮）：$L$ 增加意味着 $w/P$ \redkw{下降}；
  \item \kw{$r/P = MPK$}（粉底高亮）：$L$ 增加意味着 $r/P$ \redkw{上升}。
\end{itemize}
\end{assumpbox}

\textbf{移民均衡（第 9 页，图~\ref{fig:ch09-eq}）}：无移民成本时，劳动流向（真实）工资更高处——这里从 Foreign（初始工资低）流向 Home（初始工资高）。\kw{均衡点 B：两国工资相等}。图中：Home 工资线（MPL，向右下）与 Foreign 工资线（MPL$^*$，从右轴向左上）相交于 $B$；初始点 $A$（Home 工资 $W$）与 $A^*$（Foreign 工资 $W^*$）；移民 $\bar{L} \to$ 使 Home 劳动增加、Foreign 减少，直至共同的 $W'$。手写标注：Foreign MPL、Home MPL。

\begin{figure}[H]
  \centering
  \includegraphics[width=0.6\linewidth]{figs/ch09_fig_01.png}
  \caption{第 9 页：移民均衡——世界劳动总量固定（World amount of labor，粉底），移民使两国工资在 $B$ 点均等（$W'$）；红圈标出初始点 $A$、$A^*$ 与均衡点 $B$ 及迁移量。}
  \label{fig:ch09-eq}
\end{figure}

\textbf{总体得与失（第 10--11 页）}：于是：单一商品、L 与 K 的经济中，\term{移民使 $w/P$ 下降、$r/P$ 上升}。\redkw{Home 的原有居民总体上得益吗？K 所有者能补偿原有工人吗？}（a) Yes / b) No / c) 不确定）板书分类（高亮）：\kw{Capital Owner}（资本所有者）；\kw{Workers} $\Rightarrow$ \kw{Stayers（留下者）}、\kw{Emigrants（移出者）}、\kw{Immigrants（移入者）}；手写``(GNP/GDP)''——注意核算口径。答案：\redkw{a) Yes}（红字）——\term{得益大于损失；移民使资本所有者的得益大于原有工人的损失}。

\textbf{福利分析（第 13--15 页，图~\ref{fig:ch09-welfare}）}：劳动市场上的``消费者剩余'' $=$ \kw{K 所有者的收益}（gains of K owners，橙字）；工资线之下的矩形 $=$ \kw{工人的总工资}（total wages of workers，蓝字）。移民使劳动从 $L$ 增至 $L'$：K 所有者收益 $+(b+c)$（手写）；\kw{原有工人的总工资}下降（$-b$，手写）；手写：``\kw{Net welfare: +c}''。另注意：\term{移民工人的总工资}（灰色区域 $e$，手写标 X——不计入原有居民）。板书（第 15 页）：\kw{Home 原有居民的总收入}：\redkw{移民前 a+b+d；移民后 a+b+c+d}——\kw{总福利得益 = c}（红字大标题）。红字设问：\kw{为什么不包括 e？因为我们计算的是 Home 的原有居民}。

\begin{figure}[H]
  \centering
  \includegraphics[width=0.55\linewidth]{figs/ch09_fig_02.png}
  \caption{第 13 页：目的国（Home）的劳动市场——MPL 曲线下的面积分解为 a（K 所有者的剩余增量部分）、b、c、d（工资矩形）与 e（移民工资）；手写标注 (pre)、immigrants 与 $A$ 点。}
  \label{fig:ch09-welfare}
\end{figure}

\begin{figure}[H]
  \centering
  \includegraphics[width=0.78\linewidth]{figs/ch09_fig_08.png}
  \caption{板书（整页手写）：福利分析的微观基础——上：产品市场（$D$ = $MU$ = 支付意愿，$S$ = $MC$；消费者福利增益 $\to$ 消费者剩余；厂商增益 $\to$ 生产者剩余 $=$ 收入 $-$ 可变成本 $=$ 利润 $+$ 固定成本）；下：由 $Y = AK^{\alpha}L^{1-\alpha}$ 的 CRS 性质 $P \cdot MPL \cdot L + P \cdot MPK \cdot K = P \cdot Q$（$\pi = 0$），劳动需求 $\frac{w}{P} = \frac{\partial Y}{\partial L} = MPL$、劳动供给固定 $\bar{L}$——劳动市场上的``CS'' $\to$ 资本所有者福利、``PS'' $\to$ 工人福利（$w/P = MPL$）。}
  \label{fig:ch09-csps}
\end{figure}

\textbf{移民输出国（第 16--18 页，图~\ref{fig:ch09-emigrant}）}：失去劳动的国家呢？劳动从 $L$ 减至 $L'$（工资从 $w$ 升至 $w'$，手写 pre/post）：
\begin{itemize}
  \item \term{移民前}：K 所有者收益 \kw{a+b+c}；工人工资 \kw{e+f}（板书作答）；
  \item \term{移民后}：K 所有者 \kw{a}（自 a+b+c 减少，手写 $-(b+c)$）；\kw{留下工人} \kw{b+e}（自 e 增加，手写 $+b$）；\kw{移民者} \kw{c+d+f}（自 f 增加，手写 $+(c+d)$；移民者已不算该国居民，但板书列出其收入）；
  \item \kw{总福利得益 = d}（红字；手写 Total: $+d$）——\redkw{两国都从移民中得益}！
\end{itemize}

\begin{figure}[H]
  \centering
  \begin{minipage}{0.48\linewidth}\centering
    \includegraphics[width=\linewidth]{figs/ch09_fig_06.png}
  \end{minipage}\hfill
  \begin{minipage}{0.48\linewidth}\centering
    \includegraphics[width=\linewidth]{figs/ch09_fig_07.png}
  \end{minipage}
  \caption{移民输出国（Foreign）：左（第 17 页）——移民前的面积分解（K 所有者 a+b+c、工人 e+f）；右（第 18 页）——移民后（K 所有者 a、留下工人 b+e、移民者 c+d+f），总福利得益 $=$ d（红字）；手写标注 Capital: $-(b+c)$、Stayers: $+b$、Emigrants: $+(c+d)$ 与 stayers/immigrants 区间。}
  \label{fig:ch09-emigrant}
\end{figure}

\textbf{两国福利（第 19 页，图~\ref{fig:ch09-two}）}：一张图总结两国：\redkw{两国都得益！}Home 的得益 $=$ 三角形 C 的上半部（手写``Higher MPL''——移民在 Home 的边际产品更高），Foreign 的得益 $=$ 下半部（手写``Lower MPL$^*$''）。

\begin{figure}[H]
  \centering
  \includegraphics[width=0.55\linewidth]{figs/ch09_fig_09.png}
  \caption{第 19 页：两国福利总结——Gains to Home（C 区上部）与 Gains to Foreign（C 区下部）；手写标注 migrants 区间与 Higher MPL / Lower MPL$^*$。}
  \label{fig:ch09-two}
\end{figure}

\begin{figure}[H]
  \centering
  \includegraphics[width=0.62\linewidth]{figs/ch09_fig_20.png}
  \caption{板书（整页手写）：两国福利图（上：目的国 Home 的 a,b,c,d,e 分解；中：来源国 Foreign 的 a,b,c,d,e,f 分解）与讨论——\kw{若移民有摩擦（非无摩擦）怎么办？}即移民后 H 的工资仍高于 F！手绘 $\Delta L$ 部分迁移的图（A $\to$ B 而非到交点）：福利变化请同学自己画。}
  \label{fig:ch09-friction}
\end{figure}

\textbf{长期调整（第 20 页）}：Home 的长期调整取决于资本供给：更高的租金率可能鼓励投资或资本流入；更多 K 抬高工资、降低租金；若 $K/L$ 回到初始水平，工资也恢复——回到 Rybczynski 的长期预测。

\textbf{熟练 vs.\ 非熟练（第 21 页）}：把 K 视为\term{熟练劳动}、L 视为\term{非熟练劳动}：熟练劳动固定时，非熟练移民降低原有非熟练工资、抬高熟练工资；\redkw{对每国初始人口而言得益超过损失}（Gains exceed losses，红字）；长期调整还取决于教育与培训。

\secn{9.3 政策讨论：边境墙}

\textbf{背景（第 22 页）}：特朗普政府 2017 年主张在美墨边境建墙——阻止（非法、非熟练）移民从墨西哥进入美国。\kw{为什么美国选民支持特朗普及其提议？}非熟练国内工人短期因移民受损；尽管总体得益超过损失，\kw{却没有有效的再分配机制在运作}。\redkw{这就是故事的全部吗？}（设问）

\textbf{三个视角（第 23--25 页）}：
\begin{itemize}
  \item \term{视角 1：美国国内与移民工人的技能结构}（图~\ref{fig:ch09-skill}）：外国出生工人在\kw{最低与最高教育组}集中（手写``敲''（低技能）、``本科''标注）——中等教育水平（高中与大学毕业生）的外国出生份额小得多（10\%--15\%）；而美国本土工人约只有 10\% 处于低教育组与高教育组（手写补注）。因此：\kw{处于很低或很高教育水平的 20\% 美国本土工人倾向因移民受损；80\% 中等教育水平的本土工人倾向受益}。Giovanni Peri (UC Davis)：高技能移民工人的\kw{溢出效应}带来额外收益。
  \item \term{视角 2：国内与外国非熟练工人的``可替代性''}：若高度可替代，本土工人损失更大；若不可替代，本土工人不直接面对移民竞争。这是有争议的辩论：\kw{David Card}（UC Berkeley）认为可替代性低、损失小；\kw{George Borjas}（Harvard）认为可替代性与损失都大。
  \item \term{视角 3：``移民受害者''的分布}：若``受害者''更多分布在``摇摆州''呢？
\end{itemize}

\begin{figure}[H]
  \centering
  \includegraphics[width=0.75\linewidth]{figs/ch09_fig_10.png}
  \caption{第 23 页：美国各教育组中外国出生工人的份额——两端（0--8 年：约 70\%；博士：约 28\%）高、中间低（10\%--15\%）；手写标注低技能端与本科端；下方文字解释中间组与两端组的对照。}
  \label{fig:ch09-skill}
\end{figure}

\begin{sumbox}[移民部分小结（第 25 页原文）]
\begin{itemize}
  \item 移民的\term{方向与规模}由\redkw{跨国工资比较}直接决定（wage comparison，红字斜体）；
  \item 短期中，移民使\redkw{两国总体都得益}（to overall gains，红字斜体）：目的国\kw{资本所有者的得益}足以补偿国内\kw{工人的损失}；来源国\kw{留下者与移民者的得益}足以补偿\kw{资本所有者的损失}；
  \item 长期中资本也可以调整时，这些效应被抹平（Effects erased）；
  \item 扩展到熟练--非熟练劳动的设定。
\end{itemize}
\end{sumbox}

\secn{9.4 超越最终品贸易：外包}

\textbf{为什么需要新模型（第 27--31 页）}：此前所有贸易模型：用要素生产\term{最终品}，两国只交易\term{最终品}。但回顾第 1 讲的数据：大部分贸易是\kw{工业用品与资本品}（机械、中间投入）（图~\ref{fig:ch09-imports}）。工业用品与资本品的贸易说明\redkw{生产过程的各部分可以配置在不同地点}——\term{生产过程跨越边界的配置}。例子？（引入 Boeing 787 图~\ref{fig:ch09-boeing}：机翼——日本，尾翼——中国成都，发动机——英国/俄亥俄，座椅——法国……）；服务业的碎片化（第 32 页）：24 小时客服（Amazon 白天/印度的夜晚）。

\begin{figure}[H]
  \centering
  \includegraphics[width=0.55\linewidth]{figs/ch09_fig_11.png}
  \caption{第 31 页：Boeing 787``全球制造''示意图——机翼（日本）、中央机身（意大利 Grotaglie）、前机身（美国 Wichita）、水平尾翼（意大利 Foggia）、垂直尾翼（中国成都）、起落架（法国/英国）、发动机（英国 Rolls-Royce / 美国 GE）等。}
  \label{fig:ch09-boeing}
\end{figure}

\begin{defbox}[外包的定义（第 33 页，多色高亮原文）]
\kw{定义}：在\redkw{不同国家}提供一项服务或生产{\color{remorange}商品的各个部分}，然后{\color{defblue}在另一个地点被使用或组装成最终品}——称为\redkw{foreign sourcing（外国采购）}，或更简单地称 \redkw{offshoring（外包）}。
\begin{itemize}
  \item 外包的特征是\redkw{中间投入品的贸易}（trade in intermediate inputs）；
  \item \term{与移民的相似性}：企业能够\redkw{雇用外国工人}——尽管工人不必移动。
\end{itemize}
\end{defbox}

\textbf{今天讨论的问题（第 34 页）}：外包如何改变企业在一国进行的活动？外包对劳动市场的影响——按技能类别的劳动需求；本国与外国工人的相对工资。

\textbf{哪些活动会被外包？（第 35 页，图~\ref{fig:ch09-smile}）}：a) 取决于生产过程的\term{时间顺序}？——板书``微笑曲线（Smile Curve，手写）''：价值附加（Value Added）与技能强度（Skill intensity）在 R\&D 与营销/销售两端高、组装（Assembly）中间低；b) 还是主要取决于每项任务所需的\term{工人技能}？——手写：Developing（发展中）承接低技能端、Developed（发达）保留高技能端。

\begin{figure}[H]
  \centering
  \includegraphics[width=0.66\linewidth]{figs/ch09_fig_12.png}
  \caption{第 35 页：外包哪些活动？——上：``微笑曲线''（价值附加/技能强度 vs.\ 生产阶段：R\&D——零部件——组装——营销销售）；下：任务轴按技能排序（Low-skilled 组装 $\to$ High-skilled R\&D），手写 Developing/Developed 箭头。}
  \label{fig:ch09-smile}
\end{figure}

\textbf{企业如何从外包中获益（第 36--39 页）}：两国\term{相对工资不同}！第一条假设：\kw{Foreign 低技能劳动工资低于 Home}（高亮）：$w_L^* < w_L$；同时假设低技能劳动的\term{相对}工资在 Foreign 更低（\kw{Foreign 技能溢价更高}，高亮）：$w_L^*/w_H^* < w_L/w_H$。两条工资假设意味着：在 Foreign，低技能劳动特别便宜（\kw{绝对与相对}两个意义上都便宜，高亮）。因此\redkw{外包的收益来自 Foreign 低技能工人的更低工资}（红字）。但企业还须考虑\redkw{外包的额外成本}（costs of offshoring，红字）：建厂更高价格；通讯或运输的额外成本；决策时权衡低工资的节省与资本和贸易的额外成本；假设这些\kw{额外成本对价值链中所有活动统一适用}（apply uniformly across all the activities，红字）。

\textbf{成本记号（第 39 页）}：通讯与贸易成本概括为 \kw{$T > 1$}（高亮）；Home 高技能劳动成本 $w_H$；Home 低技能 $w_L$；通过外包的高技能成本 \kw{$T \cdot w_H^*$}（高亮）；低技能 \kw{$T \cdot w_L^*$}（高亮）。

\textbf{练习：哪种情形不可能？（第 40--41 页）}：手写注明``低技能有效工资''与``高技能有效工资''：
\begin{itemize}
  \item a) $T \cdot w_L^* < w_L$ 且 $T \cdot w_H^* < w_H$（手写：F cheap / F cheap）；
  \item b) $T \cdot w_L^* < w_L$ 且 $T \cdot w_H^* > w_H$（F cheap / H cheap）；
  \item c) $T \cdot w_L^* > w_L$ 且 $T \cdot w_H^* < w_H$——\redkw{不可能！}（板书红叉；手写注：与 $w_L^* < w_L$、$\frac{w_L^*}{w_H^*} < \frac{w_L}{w_H}$ 矛盾）；
  \item d) $T \cdot w_L^* > w_L$ 且 $T \cdot w_H^* > w_H$（H cheap / H cheap）。
\end{itemize}

\begin{figure}[H]
  \centering
  \includegraphics[width=0.6\linewidth]{figs/ch09_fig_13.png}
  \caption{板书（整页手写）：为什么 c)（$T w_L^* > w_L$ 且 $T w_H^* < w_H$）不可能——由① $T > \frac{w_L}{w_L^*}$、② $T < \frac{w_H}{w_H^*}$，联立得 $\frac{w_L}{w_L^*} < \frac{w_H}{w_H^*}$，即 $\frac{w_L}{w_H} < \frac{w_L^*}{w_H^*}$——与假设③ $\frac{w_L^*}{w_H^*} < \frac{w_L}{w_H}$ 矛盾（Contradictory to ③）。}
  \label{fig:ch09-impossible}
\end{figure}

\textbf{碎片化能够发生的唯一情形（第 42 页）}：每个国家承担正的任务份额：
\[
\kw{b)}\quad T \cdot w_L^* < w_L \ \text{且}\ T \cdot w_H^* > w_H
\]
（手写：Foreign 低技能便宜、Home 高技能便宜）——这意味着\redkw{低技能任务被外包，而多数高技能任务留在本国}（配图：任务轴上组装/零部件 vs.\ 营销/R\&D）。

\secn{9.5 外包模型的形式化：任务连续统}

\begin{assumpbox}[生产侧（第 42--43 页）]
生产一单位最终品需要许多任务``s''（高亮``skill intensity''手写），沿轴从 $s = 0$ 排到 $s = 1$——``$s$ 的连续统（\kw{continuum}，粉底高亮）''；每项任务``s''需要 ``s'' 单位高技能工人与 ``1$-$s'' 单位低技能工人（手写：H：$s\, w_H + (1-s)\, w_L$；F：$T w_H^*$、$T w_L^*$；板书示例 $s = 0.2$ 与 $0.99$ 的标注）。板书图：洋红高亮 $s$ 轴上的 $[0, 0.99]$ 区间。
\end{assumpbox}

\textbf{任务成本与最终品价格（第 43 页）}：任务 ``s'' 的成本（粉底高亮）：\kw{$C(s) = s W_H + (1-s) W_L$}；无外包时的\kw{最终品价格}（粉底高亮）：
\[
P_F = \int_{s=0}^{s=1} C(s)\, ds = \int_{s=0}^{s=1} \left[ s W_H + (1-s) W_L \right] ds .
\]

\textbf{外包决策（第 44--46 页）}：任务 ``s'' 在 Home 的成本 $C(s) = sW_H + (1-s)W_L$；外包的成本 $C^*(s) = sW_H^* + (1-s)W_L^*$（粉底标注两工资）；\kw{若 $T \cdot C^*(s) < C(s)$ 则外包}（T 粉底高亮）。
\begin{itemize}
  \item \term{情形 a}：若 $T \cdot w_L^* < w_L$ 且 $T \cdot w_H^* < w_H$（手写 $s = 0$、$s = 1$ 两端），则对所有 ``s'' 有 $T \cdot C^*(s) < C(s)$——\redkw{所有任务都被外包}（红字）；
  \item \term{情形 d}：若 $T \cdot w_L^* > w_L$ 且 $T \cdot w_H^* > w_H$，则对所有 ``s'' 有 $T \cdot C^*(s) > C(s)$——\redkw{没有任务被外包}（红字）；
  \item \term{情形 b}（图~\ref{fig:ch09-derive}板书推导）：若 $T \cdot w_L^* < w_L$ 且 $T \cdot w_H^* > w_H$：由 $sTW_H^* + (1-s)TW_L^* < sW_H + (1-s)W_L$ 整理（手写逐步：$s(TW_H^* - TW_L^*) + TW_L^* < s(W_H - W_L) + W_L$；$s\,(TW_H^* - W_H + W_L - TW_L^*) < W_L - TW_L^*$；手写标注``save cost by using cheap L''与``waste cost by using expensive H''）：
  \[
  {\color{warnred}s < \frac{W_L - T W_L^*}{(W_H - T W_H^*) + (W_L - T W_L^*)} \equiv s^*}
  \]
  \redkw{解：低于阈值 $s^*$ 的任务被外包}——\kw{低技能任务（$s < s^*$）外包；高技能任务（$s > s^*$）留在国内}。
\end{itemize}

\begin{figure}[H]
  \centering
  \includegraphics[width=0.62\linewidth]{figs/ch09_fig_13.png}
  \caption{第 46 页板书：情形 b 的阈值推导——$s^* = \frac{W_L - TW_L^*}{(W_H - TW_H^*) + (W_L - TW_L^*)}$；低技能任务（$s < s^*$）外包、高技能任务（$s > s^*$）国内生产。}
  \label{fig:ch09-derive}
\end{figure}

\textbf{图形说明（第 47 页，图~\ref{fig:ch09-graph}）}：外包成本曲线 $TC^*(s) = sTW_H^* + (1-s)TW_L^*$（黄）与国内成本线 $C(s) = sW_H + (1-s)W_L$（蓝）相交于 $s^*$：左侧\kw{offshore}（手写）、右侧\kw{stay}（手写）；手写标注两端截距 $TW_L^*$ vs.\ $W_L$、$TW_H^*$ vs.\ $W_H$。

\begin{figure}[H]
  \centering
  \includegraphics[width=0.6\linewidth]{figs/ch09_fig_14.png}
  \caption{第 47 页：外包决策的图形——两成本线在 $s^*$ 处相交；$s < s^*$ 外包（offshore）、$s > s^*$ 留守（stay）；下方任务轴标注 Assembly 到 R\&D。}
  \label{fig:ch09-graph}
\end{figure}

\secn{9.6 比较静态：贸易成本下降}

\textbf{第一组比较静态（第 48 页）}：先假设工资不变；贸易成本变化（$T$ 下降——加深一国在全球价值链中的整合，粉底标 T）。板书作答（绿勾）：\kw{T 下降 $\Rightarrow$ 更高的 $s^*$}（higher，勾）——因此\kw{更多外包}（more，勾）；\kw{更大范围的活动}（larger，勾）现在在 Foreign 进行。\redkw{Questions}（红字）：对各国熟练工人的相对需求有何影响？对各国工资不平等有何影响？

\textbf{$T$ 下降的图形（第 49--50 页）}：外包成本曲线（Upward / Downward 移动？——板书勾选 \kw{Downward}：下移；手写``\kw{+ Rotation}''——$TC^*(s) = sTW_H^* + (1-s)TW_L^* = s\,\underline{(TW_H^* - TW_L^*)} + TW_L^*$：斜率与截距同时变化，黄色新线在原橙线下方）。活动 A 与 B 之间现在在 \kw{Foreign}（板书勾选）进行（图~\ref{fig:ch09-movers}）：手写分区——``\kw{Always offshore}''（$s < s^*$）、``\kw{movers}''（$s^* < s < s^{*'}$，粉底圈出 B）、``\kw{Always stay}''（$s > s^{*'}$）。

\begin{figure}[H]
  \centering
  \begin{minipage}{0.48\linewidth}\centering
    \includegraphics[width=\linewidth]{figs/ch09_fig_15.png}
  \end{minipage}\hfill
  \begin{minipage}{0.48\linewidth}\centering
    \includegraphics[width=\linewidth]{figs/ch09_fig_16.png}
  \end{minipage}
  \caption{$T$ 下降的比较静态：左（第 49 页）——外包成本曲线下移并旋转（手写：$s\underline{(TW_H^* - TW_L^*)} + TW_L^*$），交点从 $s^*$ 右移；右（第 50 页）——新交点 $s^{*'}$，A 与 B 之间的活动（movers）转移到 Foreign；手写标注 Always offshore / movers / Always stay。}
  \label{fig:ch09-movers}
\end{figure}

\textbf{技能密集度（第 51--53 页，图~\ref{fig:ch09-skillint}）}：A 与 B 之间的活动：比 Foreign 原先做的活动（A 左侧）\kw{更技能密集}（to the left of A，蓝字）——Foreign 的平均技能密集度\kw{提高}（板书勾选 Increase）；但又比 Home 现在做 的活动（B 右侧）\kw{技能密集度更低}（to the right of B）——Home 的平均技能密集度也\kw{提高}（板书勾选 Increase）。板书手写均匀分布例子（$s^* = 0.6 \to s^{*'} = 0.8$）：$\bar{S}_{\text{Home}}$：0.3 $\to$ 0.4（↑）、$\bar{S}_{\text{Foreign}}$：0.8 $\to$ 0.9（↑）。结论（第 52 页）：\redkw{外包成本下降使两国的平均技能密集度都提高}（increase，红字）——为什么？Home 因为（新）外包的工作相对\kw{不那么}技能密集；Foreign 因为（新）外包来的工作相对\kw{更}技能密集（手写``基于`次序'``）。

\begin{figure}[H]
  \centering
  \includegraphics[width=0.62\linewidth]{figs/ch09_fig_17.png}
  \caption{第 51 页：技能密集度——A 与 B 之间的活动比 Foreign 原先做的更技能密集、比 Home 现在做的更不技能密集；手写均匀分布示例表（Before/After：$\bar S_{\text{Home}}$ 0.3→0.4，$\bar S_{\text{Foreign}}$ 0.8→0.9，均 ↑）。}
  \label{fig:ch09-skillint}
\end{figure}

\textbf{直觉（第 53 页）}：同样的任务怎么会在 Home``相对不那么密集''而在 Foreign``相对更密集''？类似例子：高中的``实验班''vs.``普通班''——假设实验班里分数最低的学生转入普通班，两个班的平均分如何变化？\kw{同一个学生在不同的班里既可以是最差的也可以是最好的}。现实例子（蓝字）：\kw{中国的落后产能转移到非洲}。

\secn{9.7 技能溢价（Skill Premium）}

\textbf{与 H-O 的对照（第 54 页）}：板书：对任意 $w_L, w_H, w_L^*, w_H^*$：$T \downarrow \Rightarrow \bar{S}_H, \bar{S}_F \uparrow$；对 \kw{H/L} 与 \kw{H$^*$/L$^*$} 的需求 $\uparrow$（手写）。比较中间品贸易（外包）与最终品贸易（Heckscher--Ohlin）：H-O 模型：若 Foreign 相对缺乏高技能劳动（在低技能任务上有 C.A.），贸易扩张\kw{降低} Foreign 对高技能劳动的相对需求、Foreign 的技能溢价\kw{下降}（``\kw{Stolper--Samuelson}''，手写；板书小例：C：$a_C H^{0.6} L^{0.4}$，S：$a_S H^{0.4} L^{0.6}$，手写数字演算）；\kw{外包}：贸易成本下降（贸易扩张）\redkw{提高}两国的平均技能密集度（increase，红字）——那外包模型中的技能溢价呢？

\textbf{Home 的技能溢价（第 55 页，图~\ref{fig:ch09-premH}）}：当 Home 的（平均）技能密集度上升时，意味着 Home 的技能溢价上升——但取决于供给形状。图：最不技能密集的活动被外包 $\Rightarrow$ Home 对熟练劳动的相对需求增加（$S_H \uparrow$，手写），相对工资上升（A $\to$ B）。手写标注轴：``high skill'' / ``low skill''。

\begin{figure}[H]
  \centering
  \begin{minipage}{0.48\linewidth}\centering
    \includegraphics[width=\linewidth]{figs/ch09_fig_18.png}
  \end{minipage}\hfill
  \begin{minipage}{0.48\linewidth}\centering
    \includegraphics[width=\linewidth]{figs/ch09_fig_19.png}
  \end{minipage}
  \caption{技能溢价：左（第 55 页）——Home：相对需求右移（$S_H \uparrow$），熟练/非熟练相对工资 $W_S/W_L$ 上升（A $\to$ B）；右（第 56 页）——Foreign：转移到 Foreign 的活动比原先的更技能密集，相对需求同样增加（$S_F \uparrow$，手写），相对工资从 $A^*$ 升至 $B^*$。}
  \label{fig:ch09-premH}
\end{figure}

\textbf{Foreign 的技能溢价（第 56 页）}：当 Foreign 的（平均）技能密集度上升时，Foreign 的技能溢价也上升——但取决于供给形状。转移给 Foreign 的活动比原先在那里做的\kw{更技能密集}；因此 Foreign 对熟练劳动的相对需求也增加；相对工资从 $A^*$ 升至 $B^*$。

\begin{thmbox}[外包模型的总预测（第 57 页）]
\term{外包模型预测：外包的扩张使\redkw{两国}的技能溢价\redkw{上升}（increase，红字）}（给定劳动相对供给无弹性）。与 H-O 模型对比：贸易扩张使（高技能丰裕的）Home 技能溢价\redkw{上升}，但使（低技能丰裕的）Foreign 技能溢价\redkw{下降}（ decrease，绿字）。\redkw{福利含义会有何不同？}（手写：$\frac{w_H}{w_L} \uparrow$ in both countries）
\end{thmbox}

\textbf{最终品价格（第 58 页）}：无外包时 $P_F = \int_{s=0}^{s=1} C(s)\,ds$；有外包时（手写：外包部分``外包''、其余``留在本国''）：
\[
P_F' = \int_{s=0}^{s=s^*} {\color{hlorange}TC^*(s)}\, ds + \int_{s=s^*}^{s=1} C(s)\, ds .
\]
由于 $s^*$ 的定义保证对所有 $s < s^*$ 有 \kw{$TC^*(s) < C(s)$}（高亮），故 $P_F' < P_F$——\kw{外包降低生产成本}，或者说，最终品出现``\kw{生产率得益}（productivity gain，粉底高亮）''。

\textbf{工人的福利（第 59--60 页）}：手写 $\frac{w_H}{w_L} \uparrow$。真实工资 $=$ 工资与价格的相对变化；\term{熟练劳动工资的涨幅总是超过价格}（手写 P $\downarrow$）；若劳动供给更\kw{无弹性}（inelastic，接近垂直），非熟练工资的\kw{降幅}超过价格；若更\kw{有弹性}（elastic，接近水平），非熟练工资降幅小于价格。理论上对非熟练劳动\redkw{模糊}：``\kw{生产率得益（更低价格）}（粉底高亮）'' vs. ``\kw{被外包替代（更低非熟练工资）}''（黄底高亮）。

\begin{figure}[H]
  \centering
  \includegraphics[width=0.55\linewidth]{figs/ch09_fig_20.png}
  \caption{板书（整页手写）：工人福利变化的论证——① 若技能溢价 $\frac{w_S}{w_L}$ 上升更多，熟练工资涨更多、非熟练工资降更多；② 外包规模增大时 $P$ 下降。Case 1（相对劳动供给有弹性）：$\Delta \frac{w_S}{w_L}$ 小（虽为负）$\Rightarrow \Delta(w_L/P) \approx \Delta w_L - \Delta P_L > 0$（非熟练真实工资上升）；Case 2（相对供给无弹性）：$\Delta \frac{w_S}{w_L}$ 大 $\Rightarrow \Delta w_L$ 大且为负 $\Rightarrow \Delta(w_L/P) \approx \Delta w_L - \Delta P_L < 0$。}
  \label{fig:ch09-cases}
\end{figure}

\secn{9.8 含义：融入全球价值链（GVC）}

\textbf{模型内的影响（第 61 页）}：手写``$T > 1$''。\kw{GVC 整合在模型中的影响}：① 更多任务被外包；② 新外包的活动比原先外包的活动更技能密集；③ 发展中国家的``技能密集度''改善；④ 发展中国家技能回报上升，从而对教育的需求改善。

\textbf{模型外的影响（第 62 页）}：\kw{GVC 整合在模型之外的影响}：学习、引进、消化（``\term{后发者优势}（latecomer's advantage）''）；改善管理与其他``软''技术；受训劳动向其他部门溢出；更容易营销产品与吸引 FDI。

\secn{9.9 模型构建方法地图}

\begin{mapbox}[移民与外包的方法地图]
\textbf{1. 单部门移民模型}（第 7--19 页）：
\begin{enumerate}
  \item 设定：一种商品、两国各自固定 $K_i$、可移民的 $L_i$；$P = 1$；$w/P = MPL$、$r/P = MPK$；
  \item \kw{移民均衡}：世界劳动条（Home $+$ Foreign）上两条 MPL 曲线相交 $\Rightarrow$ 工资均等（无摩擦时）；
  \item \kw{福利核算}（劳动市场图 $w/P = MPL$）：目的国——K 所有者 $+(b+c)$、stayers $-b$、净 $+c$；来源国——K 所有者 $-(b+c)$、stayers $+b$、emigrants $+(c+d)$、净 $+d$；\redkw{两国总福利都为正}；
  \item 扩展：移民摩擦（迁移后工资仍不等——福利三角形变小）、长期资本调整（$K/L$ 恢复 $\Rightarrow$ 回到 Rybczynski）、熟练/非熟练解释。
\end{enumerate}
\textbf{2. 外包模型}（第 33--59 页）：
\begin{enumerate}
  \item \kw{任务连续统} $s \in [0,1]$，任务 $s$ 的技能构成 $= (s, 1-s)$；成本 $C(s) = sW_H + (1-s)W_L$（国内）、$C^*(s) = sW_H^* + (1-s)W_L^*$、外包有效成本 $T \cdot C^*(s)$（$T > 1$）；
  \item \kw{工资条件}：$w_L^* < w_L$ 且 $w_L^*/w_H^* < w_L/w_H$（Foreign 低技能劳动绝对且相对便宜）$\Rightarrow$ 唯一可碎片化的情形：$T w_L^* < w_L$ 且 $T w_H^* > w_H$；
  \item \kw{外包阈值}：$s^* = \frac{W_L - TW_L^*}{(W_H - TW_H^*) + (W_L - TW_L^*)}$，$s < s^*$ 外包、$s > s^*$ 留守；
  \item \kw{比较静态}：$T \downarrow \Rightarrow$ 曲线下移$+$旋转 $\Rightarrow s^* \uparrow \Rightarrow$ movers 段转移 $\Rightarrow$ \redkw{两国平均技能密集度都上升}（次序论证/``转班''直觉）；
  \item \kw{技能溢价}：两国 $w_H/w_L$ 都上升（对比 H-O：只有丰裕国上升）；
  \item \kw{价格与福利}：$P_F' < P_F$（``生产率得益''）；非熟练工人真实工资模糊（生产率得益 vs.\ 替代效应，取决于供给弹性——Case 1/Case 2 板书）。
\end{enumerate}
\end{mapbox}

\begin{sumbox}[本讲总结：通用结论与方法]
\begin{enumerate}
  \item \term{移民}：无摩擦迁移使两国工资均等；目的国 K 所有者得益 $+$ 原有工人受损（净 $+c$），来源国 K 所有者受损 $+$ 留下者与移民者得益（净 $+d$）——\redkw{两国总福利都上升}；摩擦使迁移不完全（工资差仍存）；长期资本调整抹平工资效应（Rybczynski）。
  \item \term{技能结构 matters}：移民集中在技能分布两端——美国中等教育组（80\%）受益、两端（20\%）受损；可替代性之争（Card vs.\ Borjas）与``受害者''的政治分布解释保护主义（边境墙）。
  \item \term{外包 $=$ 中间品贸易}：企业``雇用外国工人''而工人不移动——与移民等价的生产侧延伸；Boeing 787 与 24 小时客服是生产碎片化的例子。
  \item \term{外包的收益与成本}：收益来自 Foreign 低技能工资（绝对且相对便宜：$w_L^* < w_L$ 且 $w_L^*/w_H^* < w_L/w_H$）；成本 $T > 1$ 统一作用于所有任务；唯一可碎片化的工资情形是 $T w_L^* < w_L$ 且 $T w_H^* > w_H$。
  \item \term{任务连续统与阈值}：$C(s) = sW_H + (1-s)W_L$；$s^* = \frac{W_L - TW_L^*}{(W_H - TW_H^*) + (W_L - TW_L^*)}$；$s < s^*$ 外包、$s > s^*$ 留守——低技能任务外包、高技能任务国内。
  \item \term{$T$ 下降}：外包成本线下移$+$旋转 $\Rightarrow s^*$ 上升 $\Rightarrow$ ``movers''段（介于新旧阈值之间）转移到 Foreign；\redkw{两国平均技能密集度都上升}（Home 甩掉相对低技能段、Foreign 接入相对高技能段——``转班''直觉）。
  \item \term{技能溢价}：外包使两国 $w_H/w_L$ 都上升（给定相对供给无弹性）；与 H-O（Stolper--Samuelson：只有技能丰裕国溢价上升）形成鲜明对照——这是外包模型最重要的预测差异。
  \item \term{价格与福利}：$P_F' < P_F$（生产率得益）；熟练工人确定受益（工资涨幅 $>$ 价格降幅）；非熟练工人理论上模糊——取决于相对劳动供给弹性（Case 1 弹性：真实工资升；Case 2 无弹性：真实工资降）。
  \item \term{GVC 整合的含义}：模型内——更多外包、外包任务更技能密集、发展中（承接）国技能密集度与教育回报上升；模型外——后发者优势、软技术、劳动溢出、 easier marketing 与 FDI。
\end{enumerate}
\end{sumbox}
